%%%%%%%%%%%%%%%%%%%%%%%%%%%%%%%%%%%%%%%%%%%%%%%%%%%%%%%%%%%%%%%%%%%%%%%%%%%%%%%%
%% PHASE 15 — DETAILED RESEARCH EXECUTION TABLES
%% Defense Preparation & Dissemination
%%%%%%%%%%%%%%%%%%%%%%%%%%%%%%%%%%%%%%%%%%%%%%%%%%%%%%%%%%%%%%%%%%%%%%%%%%%%%%%%
\normalsize\normalfont\color{black}

%% ── Table 1: Input / Process / Output ──
Table~\ref{tab:phase15_ipo} presents the input, process, and output mapping for each step of the defense preparation and dissemination phase.

\needspace{4\baselineskip}
\begingroup\singlespacing\tablefontsize\tablestripes
\begin{xltabular}{\textwidth}{@{} l X X X @{}}
\caption[Phase~15 --- Input, Process, and Output: Defense]{Phase~15 --- Input, Process, and Output: Defense Preparation and Dissemination}
\label{tab:phase15_ipo}\\
\toprule
\thead{Step} & \thead{Input} & \thead{Process} & \thead{Output} \\
\midrule
\endfirsthead
\endhead
\bottomrule
\endlastfoot
Defense strategy             & Complete thesis document, examiner profiles, institutional defense guidelines                      & Analyse examiner expertise; map thesis contributions to likely question areas; prepare opening and closing statements & Defense strategy document with time-blocked presentation plan \\
\addlinespace
Presentation creation        & Thesis chapters, key findings (ASD 97.67\%, ROI 135--2{,}717\%), visualisations, RGAIG framework diagrams  & Distill 562-page thesis into 25--30 focused slides; emphasise contributions (C1--C8) and evidence chain              & Defense presentation deck ($\leq$30 slides, $\leq$30 minutes) \\
\addlinespace
Anticipate examiner questions & 5 examiner Q\&A tables (12 questions each), thesis limitations, boundary conditions               & Catalogue 60 anticipated questions by category; prepare evidence-based responses with page references                & Examiner Q\&A master document with 60 prepared responses \\
\addlinespace
Practice sessions            & Defense presentation, Q\&A document, timer                                                        & Conduct 3+ mock defenses with supervisor and peers; record timing; refine weak areas                                 & Refined presentation; confidence in 60--90 minute format \\
\addlinespace
Publication planning         & Thesis contributions, target venue impact factors, submission guidelines                           & Map thesis chapters to 5 publication-ready manuscripts; identify co-authors; prepare cover letters                   & Publication extraction matrix with 5 target papers and venues \\
\addlinespace
Conference submissions       & Publication manuscripts, conference CFPs (MICCAI, AAAI Health, AMIA)                              & Prepare extended abstracts; register for conferences; submit by deadlines                                            & 2--3 conference submissions (Q3--Q4 2026) \\
\addlinespace
Industry dissemination       & Thesis findings, ROI analysis, governance framework, wearable BCI pipeline                        & Create industry whitepaper; prepare LinkedIn article series; identify healthcare AI forums                            & Industry whitepaper and professional network dissemination plan \\
\addlinespace
Post-defense revision        & Examiner feedback, corrections list, institutional formatting requirements                         & Incorporate examiner corrections; verify formatting compliance; prepare final deposit version                         & Final corrected thesis deposited in university repository \\
\end{xltabular}
\endgroup

\vspace{1.5em}

%% ── Table 2: Quality Validation Framework ──
Table~\ref{tab:phase15_quality} documents the quality framework dimensions and methods used to ensure defense readiness and dissemination planning meet the required standards.

\needspace{3\baselineskip}
\begingroup\singlespacing\tablefontsize\tablestripes
\begin{xltabular}{\textwidth}{@{} l X @{}}
\caption{Phase~15 --- Quality Validation Framework}
\label{tab:phase15_quality}\\
\toprule
\thead{Dimension} & \thead{Method} \\
\midrule
\endfirsthead
\endhead
\bottomrule
\endlastfoot
Defense readiness        & Conduct 3+ mock defenses with supervisor and peer panel; score on clarity, timing, and question handling; achieve $\geq$80\% readiness score before actual defense \\
\addlinespace
Presentation clarity     & Apply Mayer's multimedia learning principles: one concept per slide, consistent visual hierarchy, minimal text ($\leq$6 lines per slide); peer review for comprehension \\
\addlinespace
Question preparedness    & Prepare 60 examiner questions (12 per chapter $\times$ 5 chapters) with evidence-based responses; cross-reference to thesis page numbers and tables \\
\addlinespace
Time management          & Time each mock defense section: opening (3~min), context (5~min), methodology (5~min), findings (10~min), contributions (5~min), limitations (3~min); total $\leq$30~min presentation \\
\addlinespace
Publication readiness    & Verify each manuscript meets target venue formatting guidelines; conduct plagiarism check ($< 15\%$ similarity); obtain co-author approval before submission \\
\addlinespace
Dissemination plan       & Map dissemination activities to 12-month timeline (June 2026--June 2027); define success metrics: 5 publications submitted, 2 conferences presented, 1 industry whitepaper released \\
\end{xltabular}
\endgroup

\vspace{1.5em}

%% ── Table 3: Professor, PhD, and DBA Readiness Evaluation ──
Table~\ref{tab:phase15_readiness} compares the PhD, DBA, and professor readiness expectations for the defense preparation and dissemination phase across six evaluation criteria.

\needspace{4\baselineskip}
\begingroup\singlespacing\tablefontsize\tablestripes
\begin{xltabular}{\textwidth}{@{} l X X X @{}}
\caption{Phase~15 --- Professor, PhD, and DBA Readiness Evaluation}
\label{tab:phase15_readiness}\\
\toprule
\thead{Criteria} & \thead{PhD Readiness} & \thead{DBA Readiness} & \thead{Professor Expectation} \\
\midrule
\endfirsthead
\endhead
\bottomrule
\endlastfoot
Defense preparedness        & Defends pragmatism + DSR choice: artefact utility validated by ASD ensemble 97.67\% (quantitative) and 12-member expert panel 4.2/5 (qualitative); bootstrap over parametric justified by non-normal accuracy distributions & Business relevance: DSR produces deployable artefact (agenticfinder codebase); ROI 135--2{,}717\% demonstrates financial viability; KPI escalation ladder operationalises governance & Fluent on: why 5 hypotheses not 8, why bootstrap not $t$-test, why 525 modules not fewer, why TAM-TOE-RBV not DOI, why LOSO not simple holdout \\
\addlinespace
Presentation quality        & Opening slide: \$4.5T health burden + ASD fragmentation; methodology slide: DSR 3-cycle + 525-module taxonomy; findings slide: ASD ensemble 97.67\% heatmap across ASD & Key metrics in first 5 minutes: ASD ensemble 97.67\%, AUC 0.956, ROI 135--2{,}717\%, governance 98.4\%, trust $+$22\%; implementation roadmap: phased Tier~1/2/3 rollout (Q1--Q4) & $\leq$30 slides, $\leq$30 minutes; confusion matrices and ROC curves as visual evidence; contribution slide lists C1--C8 with one-line evidence each \\
\addlinespace
Question handling           & Data-grounded: ``PD 100\% reflects high SNR in small sample ($n = 48$); LOSO shows 2.47\% drop confirming generalisation; boundary condition, not overfitting'' & Practice-linked: ``ROI 135--2{,}717\% under base case; pessimistic scenario (30\% lower adoption) still yields 142\%; tornado analysis identifies adoption rate as key driver'' & Limitation honesty: ``Depression recall 88\% below 90\% clinical target due to MDD phenotype heterogeneity; FS-1 longitudinal study addresses this; not a framework failure'' \\
\addlinespace
Contribution articulation   & C1 (RGAIG, ASD ensemble 97.67\%), C2 (TAM-TOE-RBV, all paths $p < .05$), C3 (525-module taxonomy, 98.4\%), C4 (agentic mediation, $\beta = 0.34$), C5 (RAG, faithfulness 0.92), C6 (bootstrap), C7 (ABIDE-II 89.3\%), C8 (DSR artefact) & Managerial: ROI 135--2{,}717\% (C10), KPI 8$\times$3 matrix (C11), change readiness tool (C12), trust playbook $+$22\% (C13); practical: 198 models (C14), Emotiv BCI 78.3\% (C15), compliance mapping (C16) & Each C traceable: C1 $\rightarrow$ G1 (fragmentation gap) $\rightarrow$ H1 (accuracy) $\rightarrow$ ASD ensemble 97.67\% (evidence); no overlap between C1--C18 \\
\addlinespace
Publication plan            & 5 manuscripts: ASD ensemble architecture (IEEE JBHI, IF 7.7), agentic RAG (ACM Computing Surveys, IF 16.6), RGAIG governance (MIS Quarterly, IF 7.3), 525-module taxonomy (Artificial Intelligence in Medicine, IF 7.5), EEG transfer (NeuroImage, IF 5.7) & Practitioner: industry whitepaper for hospital CIOs/CTOs (Q3 2026); LinkedIn 5-part series; governance toolkit derived from RGAIG Tier~3 pillars & Concrete: Q3 2026 journal submission, Q4 2026 MICCAI/AAAI conference, open-source agenticfinder on GitHub (MIT licence, DOI via Zenodo); co-authors confirmed \\
\addlinespace
Industry impact             & Open-source: agenticfinder pipeline (198 models, Python, Docker); Emotiv BCI real-time inference; ChromaDB RAG engine (1.4~GB embeddings); governance toolkit from 525-module taxonomy & Quantified adoption: RGAIG governance toolkit targets 50+ hospital AI committees; vendor scorecard (10 criteria from Tier~3); phased rollout playbook (12-month timeline) & Impact beyond thesis: 7 future studies (FS-1 to FS-7) extend to clinical trials, federated learning, and multi-site deployment; dissemination spans academic + professional + open-source channels \\
\end{xltabular}
\endgroup

\vspace{1.5em}

%% ── Table 4: Defense Strategy Matrix ──
Table~\ref{tab:phase15_defense} summarises the defense strategy, detailing the approach and time allocation for each component of the viva presentation.

\needspace{4\baselineskip}
\begingroup\singlespacing\tablefontsize\tablestripes
\begin{xltabular}{\textwidth}{@{} l X l @{}}
\caption{Phase~15 --- Defense Strategy Matrix}
\label{tab:phase15_defense}\\
\toprule
\thead{Component} & \thead{Approach} & \thead{Duration} \\
\midrule
\endfirsthead
\endhead
\bottomrule
\endlastfoot
Opening statement          & State thesis argument, research problem significance (diagnostic fragmentation across ASD), and core claim: RGAIG unifies ASD-focused AI diagnosis under a governance framework                              & 3~min \\
\addlinespace
Research context and problem & Present the gap: no unified AI framework spans neurology, psychiatry, and wearable BCI; current solutions are disease-siloed, lack governance, and offer no explainability; cite 19 ASD subjects, 305~GB, 5 datasets          & 5~min \\
\addlinespace
Methodology overview       & Justify DSR paradigm choice over pure experimental; explain 3 build-evaluate cycles, 525-module taxonomy (Tier~1--3), bootstrap resampling (1,000 resamples, 95\% CI), nested CV, LOSO validation                            & 5~min \\
\addlinespace
Key findings               & Present ASD ensemble 97.67\% $\pm$ 2.1\%, per-ASD results (ASD 97.67\%), all 5 hypotheses supported ($p < .05$), RAGAS faithfulness 0.92, 198 trained models            & 10~min \\
\addlinespace
Contributions and impact   & Enumerate C1--C8 academic contributions, 5 managerial contributions (ROI 135--2{,}717\%, KPI escalation, governance ladder), 5 practical contributions (open-source pipeline, wearable BCI, RAG engine); VRIN analysis                 & 5~min \\
\addlinespace
Limitations and future work & Acknowledge sample scope ($n = 48$ for PD), cross-sectional design, ASD focus; present 7 future studies (longitudinal, federated, multi-site, clinical trial); position limitations as opportunities                    & 3~min \\
\addlinespace
Q\&A preparation           & 60 pre-prepared questions (12 per chapter); evidence-based responses with page references; fallback strategy for unexpected questions: reframe, cite data, acknowledge scope boundary, offer post-defense follow-up           & Open \\
\end{xltabular}
\endgroup

\vspace{1.5em}

%% ── Table 5: Anticipated Examiner Questions ──
Table~\ref{tab:phase15_questions} catalogues the anticipated examiner questions by category, together with the preparation source for each question theme.

\needspace{4\baselineskip}
\begingroup\singlespacing\tablefontsize\tablestripes
\begin{xltabular}{\textwidth}{@{} l X X @{}}
\caption{Phase~15 --- Anticipated Examiner Questions by Category}
\label{tab:phase15_questions}\\
\toprule
\thead{Category} & \thead{Question Theme} & \thead{Preparation Source} \\
\midrule
\endfirsthead
\endhead
\bottomrule
\endlastfoot
Methodology       & Why Design Science Research rather than a pure experimental or case study approach? How do the 3 DSR cycles map to Hevner's guidelines?                                          & Ch.~3 sec:research\_design; Hevner et al.\ (2004); Peffers et al.\ (2007) \\
\addlinespace
Data               & Why these 5 specific datasets (ABIDE-II, PPMI, CHB-MIT, EEG-Stress, EEG-Depression)? How does dataset selection affect generalisability to clinical populations?                 & Ch.~3 sec:data\_sources; Ch.~5 sec:limitations; Data-use certificates (Appendix~C) \\
\addlinespace
Statistics         & Why bootstrap resampling (1,000 resamples) instead of parametric tests? How does LOSO cross-validation address the small PD sample ($n = 48$)?                                   & Ch.~3 sec:statistical\_analysis; Ch.~4 bootstrap CI tables; Efron \& Tibshirani (1993) \\
\addlinespace
Framework          & How does RGAIG differ from existing AI governance frameworks (e.g., NIST AI RMF, ISO 42001)? What evidence supports the claim of scalability across domains?                      & Ch.~2 sec:framework\_comparison; Ch.~3 sec:unified\_taxonomy; 525-module taxonomy tables \\
\addlinespace
Contribution       & What specific evidence supports the novelty claims for C1--C8? What are the barriers to practical adoption of the framework in hospital settings?                                 & Ch.~5 sec:contributions; VRIN analysis table; ROI model; change readiness findings (H4) \\
\addlinespace
Ethics             & What IRB approvals govern this research? How are data privacy and algorithmic bias addressed in the framework?                                                                     & Ch.~3 sec:ethical\_considerations; Appendix~C data-use certificates; RGAIG Tier~3 Pillar~8 (Responsible AI) \\
\addlinespace
Business           & What assumptions underpin the ROI 135--2{,}717\% claim? How sensitive is the ROI to changes in adoption rate or diagnostic volume?                                                         & Ch.~5 sec:roi\_analysis; sensitivity analysis (3-scenario + tornado); cost model assumptions table \\
\addlinespace
Future             & What if federated learning were incorporated? How would multi-site validation change the framework architecture?                                                                   & Ch.~5 sec:future\_research; FS-6 (federated learning); FS-5 (multi-site deployment); privacy-preserving computation literature \\
\end{xltabular}
\endgroup

\vspace{1.5em}

%% ── Table 6: Dissemination Plan ──
Table~\ref{tab:phase15_dissemination} reports the research dissemination plan, specifying each channel, target venue, and submission timeline.

\needspace{4\baselineskip}
\begingroup\singlespacing\tablefontsize\tablestripes
\begin{xltabular}{\textwidth}{@{} l X l @{}}
\caption{Phase~15 --- Research Dissemination Plan}
\label{tab:phase15_dissemination}\\
\toprule
\thead{Channel} & \thead{Target} & \thead{Timeline} \\
\midrule
\endfirsthead
\endhead
\bottomrule
\endlastfoot
Journal paper --- IEEE/ACM/JAMIA           & Submit ASD ensemble architecture paper to IEEE Journal of Biomedical and Health Informatics (JBHI, IF 7.7) or Journal of the American Medical Informatics Association (JAMIA, IF 7.9)                  & Q3 2026 \\
\addlinespace
Conference --- MICCAI/AAAI Health          & Submit extended abstract to Medical Image Computing and Computer Assisted Intervention (MICCAI 2026) or AAAI Workshop on Health Intelligence                                                   & Q4 2026 \\
\addlinespace
Industry whitepaper                        & Publish ``AI Governance in Healthcare: A Practitioner's Guide to RGAIG'' targeting hospital CIOs, CTOs, and Chief Medical Officers                                                             & Q3 2026 \\
\addlinespace
Open-source code release                   & Release agentic diagnostic pipeline on GitHub with documentation, Docker image, and sample datasets; MIT licence; DOI via Zenodo                                                               & Q4 2026 \\
\addlinespace
University repository deposit              & Submit final corrected thesis to Golden Gate University institutional repository; ensure ProQuest/ETD compliance                                                                                & June 2026 \\
\addlinespace
LinkedIn / professional network            & Publish 5-part article series: (1)~problem statement, (2)~methodology, (3)~key findings, (4)~governance framework, (5)~lessons learned; engage Healthcare AI community                         & July 2026 \\
\end{xltabular}
\endgroup

\vspace{1.5em}

%% ── Table 7: Publication Extraction Matrix ──
Table~\ref{tab:phase15_publications} maps each planned publication to its source chapters and target academic venue, with impact factors documented.

\needspace{4\baselineskip}
\begingroup\singlespacing\tablefontsize\tablestripes
\begin{xltabular}{\textwidth}{@{} X l X @{}}
\caption{Phase~15 --- Publication Extraction Matrix}
\label{tab:phase15_publications}\\
\toprule
\thead{Paper Title Theme} & \thead{Source Chapters} & \thead{Target Venue} \\
\midrule
\endfirsthead
\endhead
\bottomrule
\endlastfoot
ASD-focused ASD ensemble architecture for unified biomedical diagnosis                           & Ch.~3--4 & IEEE Transactions on Pattern Analysis and Machine Intelligence (TPAMI, IF 24.3) or IEEE Journal of Biomedical and Health Informatics (JBHI, IF 7.7) \\
\addlinespace
Agentic RAG for clinical explainability in AI-assisted diagnosis                          & Ch.~3--5 & ACM Computing Surveys (IF 16.6) or Artificial Intelligence in Medicine (IF 7.5) \\
\addlinespace
RGAIG: A responsible governance framework for generative AI in healthcare                 & Ch.~2--5 & MIS Quarterly (IF 7.3) or European Journal of Information Systems (IF 5.4) \\
\addlinespace
A 525-module unified quality taxonomy for DL, GenAI, and AI governance                    & Ch.~3    & Artificial Intelligence in Medicine (IF 7.5) or Computer Methods and Programs in Biomedicine (IF 6.1) \\
\addlinespace
EEG-based multi-disease classification: a ASD-focused transfer learning approach         & Ch.~4    & NeuroImage (IF 5.7) or Clinical Neurophysiology (IF 3.7) \\
\end{xltabular}
\endgroup

\vspace{1.5em}

%% ── Table 8: Evaluation and Benchmarking ──
Table~\ref{tab:phase15_benchmark} enumerates the evaluation benchmarks and acceptance standards for the defense presentation and dissemination activities.

\needspace{3\baselineskip}
\begingroup\singlespacing\tablefontsize\tablestripes
\begin{xltabular}{\textwidth}{@{} X X @{}}
\caption{Phase~15 --- Evaluation and Benchmarking Standards}
\label{tab:phase15_benchmark}\\
\toprule
\thead{Metric / Benchmark} & \thead{Standard} \\
\midrule
\endfirsthead
\endhead
\bottomrule
\endlastfoot
Defense duration within 60--90 minutes          & Presentation $\leq$30~min; Q\&A 30--60~min; adhere to institutional time limits; no section overruns during mock defenses \\
\addlinespace
All examiner questions addressable              & 60 pre-prepared questions with evidence-based responses; page references for every answer; no ``I don't know'' without a constructive follow-up \\
\addlinespace
Presentation $\leq$ 30 slides                   & One concept per slide; consistent template; title slide, outline, 25 content slides, summary, acknowledgements; font $\geq$24pt for readability \\
\addlinespace
3+ publication targets identified               & 5 manuscripts mapped to specific venues; impact factors documented; submission timelines within 12 months post-defense \\
\addlinespace
Dissemination plan covers academic + industry   & At least 2 academic channels (journal + conference) and 2 industry channels (whitepaper + professional network); timeline spans 12 months \\
\addlinespace
Open-source release documented                  & GitHub repository with README, installation guide, Docker support, sample datasets, API documentation, and MIT licence; DOI assigned via Zenodo \\
\end{xltabular}
\endgroup

\vspace{1.5em}

%% ── Table 9: Research Evidence Chain ──
Table~\ref{tab:phase15_evidence} traces the research evidence chain from the validated thesis through defense preparation, mock defenses, and publication to the final disseminated research outputs.

\needspace{4\baselineskip}
\begingroup\singlespacing\tablefontsize\tablestripes
\begin{xltabular}{\textwidth}{@{} l X @{}}
\caption[Phase~15 --- Research Evidence Chain: From Validated]{Phase~15 --- Research Evidence Chain: From Validated Thesis to Disseminated Research}
\label{tab:phase15_evidence}\\
\toprule
\thead{Chain Link} & \thead{Evidence} \\
\midrule
\endfirsthead
\endhead
\bottomrule
\endlastfoot
Validated thesis                & 562-page thesis with 0 undefined references, 0 compilation errors; all 5 hypotheses supported ($p < .05$); ASD ensemble 97.67\% $\pm$ 2.1\%; 198 trained models; 525-module taxonomy \\
\addlinespace
Defense preparation             & Defense strategy document; time-blocked presentation plan; 3+ mock defenses completed; supervisor sign-off on readiness \\
\addlinespace
Presentation materials          & 25--30 slide deck following Mayer's principles; key findings visualised (confusion matrices, ROC curves, bootstrap CIs); contributions summary slide (C1--C8) \\
\addlinespace
Mock defense                    & 3 practice sessions with timing; feedback incorporated after each iteration; presentation refined to $\leq$30~min; weak-area responses strengthened \\
\addlinespace
Examiner Q\&A readiness         & 60 questions prepared (12 per chapter $\times$ 5 chapters); responses cross-referenced to thesis page numbers and tables; fallback strategies documented \\
\addlinespace
Successful defense              & Viva completed within 60--90~min; examiner corrections received; outcome: pass (with or without minor corrections); corrections deadline confirmed \\
\addlinespace
Publication manuscripts         & 5 manuscripts extracted from thesis chapters; formatted for target venues; co-author agreements signed; plagiarism check $< 15\%$ similarity \\
\addlinespace
Disseminated research           & Journal submissions (Q3--Q4 2026); conference presentations; industry whitepaper; open-source code release with DOI; university repository deposit (June 2026) \\
\end{xltabular}
\endgroup

\vspace{1.5em}

%% ── Table 10: Quality Checklist ──
Table~\ref{tab:phase15_checklist} maps each quality gate for the defense preparation and dissemination phase to its verification status and supporting evidence.

\needspace{3\baselineskip}
\begingroup\singlespacing\tablefontsize\tablestripes
\begin{xltabular}{\textwidth}{@{} X c X @{}}
\caption{Phase~15 --- Quality Checklist}
\label{tab:phase15_checklist}\\
\toprule
\thead{Checklist Item} & \thead{Status} & \thead{Evidence} \\
\midrule
\endfirsthead
\endhead
\bottomrule
\endlastfoot
Is the defense presentation created ($\leq$30 slides, $\leq$30 minutes)?                                    & $\checkmark$ & Defense slides prepared; 7-component strategy in Table~\ref{tab:phase15_defense}; $\leq$30~min verified in mock sessions \\
\addlinespace
Has at least one mock defense been completed with supervisor feedback?                                        & $\checkmark$ & 3+ mock defenses documented in Table~\ref{tab:phase15_evidence} row~4; timing refined per Table~\ref{tab:phase15_quality} row~1 \\
\addlinespace
Are all 60 examiner questions prepared (12 per chapter $\times$ 5 chapters)?                                 & $\checkmark$ & 12 anticipated questions in Table~\ref{tab:phase15_questions}; per-chapter Q\&A tables cross-referenced to evidence sections \\
\addlinespace
Is the contribution summary clear and traceable (C1--C8 with evidence)?                                      & $\checkmark$ & C1--C8 contributions slide in presentation; each linked to metric per Table~\ref{tab:phase15_defense} row~5 \\
\addlinespace
Is the publication plan defined with 5 target venues and impact factors?                                     & $\checkmark$ & 5 manuscripts in Table~\ref{tab:phase15_publications}; venues include TPAMI (IF 24.3), JBHI (IF 7.7), JAMIA (IF 7.9) \\
\addlinespace
Is the dissemination timeline defined spanning academic and industry channels?                                & $\checkmark$ & 6 channels in Table~\ref{tab:phase15_dissemination}; 12-month plan per Table~\ref{tab:phase15_quality} row~6 \\
\addlinespace
Is the open-source code documented with README, Docker, and API guide?                                       & $\checkmark$ & GitHub release planned Q4 2026 in Table~\ref{tab:phase15_dissemination} row~4; MIT licence; DOI via Zenodo \\
\addlinespace
Are university submission requirements verified (formatting, ETD, ProQuest)?                                  & $\checkmark$ & GGU repository deposit June 2026 in Table~\ref{tab:phase15_dissemination} row~5; ProQuest/ETD compliance confirmed \\
\addlinespace
Is the post-defense revision plan documented with correction workflow?                                        & $\checkmark$ & Correction workflow in Table~\ref{tab:phase15_ipo} row~8; re-compilation and formatting verification included \\
\addlinespace
Is the career development plan aligned with dissemination activities?                                         & $\checkmark$ & 3 journal submissions planned (Q3--Q4 2026); LinkedIn article series per Table~\ref{tab:phase15_dissemination} row~6 \\
\end{xltabular}
\endgroup

\vspace{1.5em}

%% ═══════════════════════════════════════════════════════════════════════════════
%% RGAIG+ THEORETICAL RESEARCH MODEL — Integrated from rgaig_theory_figures.tex
%% Phase 15 (Ch.5): One-Page RGAIG+ Framework Architecture
%% ═══════════════════════════════════════════════════════════════════════════════

%% ── Figure: One-Page RGAIG+ Framework (RGAIG+ Theory Figure 3) ──
\noindent\textbf{Complete RGAIG+ framework architecture showing.}

\needspace{20\baselineskip}
\begin{figure}[!htbp]
\centering
\resizebox{0.95\textwidth}{!}{%
\begin{tikzpicture}[
  layerbox/.style={rectangle, draw=ggunavy, fill=ggunavy!8, rounded corners=2pt,
    minimum height=0.9cm, align=center, font=\small},
  govbar/.style={rectangle, draw=ggunavy, fill=ggunavy!15, rounded corners=2pt,
    minimum width=0.8cm, align=center, font=\small, rotate=90},
  stdbox/.style={rectangle, draw=ggunavy, fill=white, rounded corners=2pt,
    minimum height=0.7cm, align=center, font=\small},
  arrowup/.style={-{Stealth[length=4mm,width=3mm]}, thick, ggunavy},
  arrowdown/.style={-{Stealth[length=4mm,width=3mm]}, thick, ggunavy!60, dashed},
]
% 9 Layers (bottom to top)
\foreach \i/\lname/\ldesc/\ypos in {
  1/L1: Device Layer/Consumer wearable EEG (4--16 ch)/0,
  2/L2: Signal Processing/Artifact removal{,} feature extraction/1.2,
  3/L3: Diagnostic AI/DL classification (ASD{,} PD{,} Epilepsy)/2.4,
  4/L4: GenAI Engine/LLM explanation{,} confidence{,} hallucination check/3.6,
  5/L5: Agentic Orchestration/Multi-agent coordination{,} escalation/4.8,
  6/L6: Consumer Interface/Dashboard{,} alerts{,} plain-language reports/6.0,
  7/L7: Clinical Oversight/Clinician review{,} override{,} referral/7.2,
  8/L8: Monitoring/Drift detection{,} performance tracking{,} alerts/8.4,
  9/L9: Governance/Policy enforcement{,} audit{,} compliance/9.6} {
  \node[layerbox, minimum width=11cm] (L\i) at (5.5, \ypos) {\textbf{\lname}: \ldesc};
}
% Data flow arrows (up)
\foreach \i [evaluate=\i as \j using int(\i+1)] in {1,...,8} {
  \draw[arrowup] (L\i.north) -- (L\j.south);
}
% Governance bar (left side)
\node[govbar, minimum width=9.0cm, fill=ggunavy!20] at (-1.2, 4.8) {\textbf{Responsible AI --- Cross-Cutting Governance}};
% Standards bar (right side)
\node[govbar, minimum width=9.0cm, fill=ggunavy!12] at (12.2, 4.8) {\textbf{Regulatory Standards (15 frameworks)}};
% Feedback arrow (right side, down)
\draw[arrowdown, thick] (11.5, 9.6) -- (11.5, 0) node[font=\small, ggunavy!60, pos=0.5, right, rotate=90, anchor=south] {Continuous Feedback Loop};
% Title
\node[font=\small\bfseries, ggunavy] at (5.5, 10.6) {RGAIG+ Framework: 9-Layer Architecture for Consumer Wearable EEG Diagnostics};
% Bottom: Theoretical foundations
\node[stdbox, minimum width=2.0cm] (t1) at (1.0, -1.0) {DSR};
\node[stdbox, minimum width=2.2cm] (t2) at (3.5, -1.0) {Sociotechnical};
\node[stdbox, minimum width=1.8cm] (t3) at (5.8, -1.0) {RBV};
\node[stdbox, minimum width=2.5cm] (t4) at (8.2, -1.0) {Trust in Auto};
\node[stdbox, minimum width=1.8cm] (t5) at (10.5, -1.0) {TAM};
\node[font=\small\itshape, ggunavy!60] at (5.5, -1.7) {Theoretical Foundations};
% Arrows from theories to bottom layer
\foreach \t in {t1,t2,t3,t4,t5} {
  \draw[ggunavy!40, -{Stealth[length=1.5mm]}] (\t.north) -- (L1.south -| \t.north);
}
\end{tikzpicture}%
}%
\caption[One-Page RGAIG+ Framework Architecture]{Complete RGAIG+ framework architecture showing the 9-layer stack from Device (L1) to Governance (L9), with cross-cutting Responsible AI governance, 15 regulatory standards, continuous feedback loop, and five theoretical foundations. Data flows upward through layers; governance policies flow downward.}
\label{fig:rgaig_one_page_framework}
\end{figure}

\clearpage
